% Reviewed translation: PDF pages 261--264 / printed pages 247--250.
% The original scan is authoritative; mathematical tokens were checked against it.
\MathBookSourcePage{261}
除 \eqref{eq:incompressible-hhj-restriction-orthogonality} 外,
双线性形式 $b_h$ 还具有下列显著性质.

\begin{Lemma}\label{lem:incompressible-hhj-kernel-equivalence}
对每个 $\boldsymbol\tau_h\in\mathcal X_h$, 有等价关系
\begin{equation}\label{eq:incompressible-hhj-kernel-equivalence}
\{b_h(\boldsymbol\tau_h,\phi_h)=0
\ \forall\phi_h\in\Phi_h\}
\quad\Longleftrightarrow\quad
\{b_h(\boldsymbol\tau_h,\phi)=0
\ \forall\phi\in\widetilde\Psi\cap C^0(\overline\Omega)\}.
\end{equation}
\end{Lemma}

\begin{Proof}
显然只需证明由左向右的方向.
为此, 取 $\phi\in C^0(\overline\Omega)$,
$\phi|_\kappa\in H^2(\kappa)$, 因而 $\phi\in H^1(\Omega)$,
并证明
\begin{equation}\label{eq:incompressible-hhj-interpolation-orthogonality}
b_h(\boldsymbol\tau_h,\phi-I_h\phi)=0
\qquad\forall\boldsymbol\tau_h\in\mathcal X_h,
\end{equation}
其中 $I_h$ 是附录中定义的插值算子.
两次分部积分给出
\[
\begin{split}
\int_\kappa
\frac{\partial^2\phi}{\partial x_i\partial x_j}f\,dx
={}&\int_\kappa
\frac{\partial^2f}{\partial x_i\partial x_j}\phi\,dx
-\int_{\partial\kappa}
\frac{\partial f}{\partial x_i}\phi n_j\,ds\\
&+\int_{\partial\kappa}
f\frac{\partial\phi}{\partial x_i}n_j\,ds
\end{split}
\]
对所有 $f,\phi\in H^2(\kappa)$ 成立.
当 $f\in P_{l-1}$ 时, 插值算子的矩条件给出
\[
\int_\kappa f
\frac{\partial^2(\phi-I_h\phi)}{\partial x_i\partial x_j}\,dx
=\int_{\partial\kappa}fn_j
\frac{\partial(\phi-I_h\phi)}{\partial x_i}\,ds.
\]
代入 $b_h$ 的定义, 得
\[
b_h(\boldsymbol\tau_h,\phi-I_h\phi)
=-\sum_{\kappa\in\mathcal T_h}
\int_{\partial\kappa}M_{nt}(\boldsymbol\tau_h)
\frac{\partial(\phi-I_h\phi)}{\partial t}\,ds.
\]
再在每条边上分部积分并再次应用插值矩条件,
即得 $b_h(\boldsymbol\tau_h,\phi-I_h\phi)=0$.
这证明 \eqref{eq:incompressible-hhj-kernel-equivalence}.
\end{Proof}

由插值算子的定义容易推出,
Lemma~\ref{lem:incompressible-sfvp-mesh-two-approximation}
对 $\phi_h=I_hv$ 仍成立:
\begin{equation}\label{eq:incompressible-hhj-stream-interpolation-error}
\|v-I_hv\|_{2,h}
\leq Ch^{k-2}|v|_{k,\Omega}
\qquad\forall v\in H^k(\Omega),
\end{equation}
其中假设 $\mathcal T_h$ 正则.
结合 \eqref{eq:incompressible-hhj-interpolation-orthogonality}
和 \eqref{eq:incompressible-hhj-stream-interpolation-error}, 得

\begin{Corollary}\label{cor:incompressible-hhj-stream-interpolant}
附录定义的算子 $I_h$ 满足
\[
b_h(\boldsymbol\tau_h,\phi-I_h\phi)=0
\]
对所有分片 $H^2$ 的
$\phi\in C^0(\overline\Omega)$
和所有 $\boldsymbol\tau_h\in\mathcal X_h$ 成立.
此外, 若 $\mathcal T_h$ 正则,
则存在与 $h$ 和 $\phi$ 无关的常数 $C>0$, 使
\[
\|\phi-I_h\phi\|_{2,h}
\leq Ch^{k-2}|\phi|_{k,\Omega}
\qquad\forall\phi\in H^k(\Omega),
\]
其中实数 $k\in[2,l+1]$.
\end{Corollary}

现在转向 inf--sup 条件.
先把自己限制在张量空间
$\mathcal X\cap H^1(\Omega)^4$ 上,
特别考虑形如
\[
\boldsymbol\tau
=\begin{pmatrix}
-\partial\theta/\partial x_2&\partial\theta/\partial x_1\\
\partial\theta/\partial x_1&\partial\theta/\partial x_2
\end{pmatrix}
\]
的张量. 所有这些张量满足
\[
b_h(\boldsymbol\tau,\phi_h)
=(\operatorname{curl}\theta,\operatorname{curl}\phi_h)
\qquad\forall\phi_h\in\Phi_h,
\]
并且
\[
\|\boldsymbol\tau\|_{0,h}
\leq\sqrt2\left(
\|\theta\|_{0,\Omega}^2
+h\|\theta\|_{0,\Gamma_h}^2
\right)^{1/2}.
\]

\MathBookSourcePage{262}
应用 Lemma~\ref{lem:incompressible-sfvp-mesh-inf-sup}, 得初步结果
\[
\sup_{\boldsymbol\tau\in\mathcal X\cap H^1(\Omega)^4}
\frac{b_h(\boldsymbol\tau,\phi_h)}
{\|\boldsymbol\tau\|_{0,h}}
\geq\frac1{\sqrt2}\beta^*\|\phi_h\|_{2,h}
\qquad\forall\phi_h\in\Phi_h,
\]
其中 $\beta^*$ 是
Lemma~\ref{lem:incompressible-sfvp-mesh-inf-sup} 中的常数,
并假设 $\Omega$ 为凸域且剖分一致正则.
由 \eqref{eq:incompressible-hhj-restriction-orthogonality}
和 \eqref{eq:incompressible-hhj-restriction-stability},
它蕴含 $\mathcal X_h$ 上的 inf--sup 条件.

\begin{Lemma}\label{lem:incompressible-hhj-inf-sup}
令 $\Omega$ 是有界凸多边形,
$\mathcal T_h$ 是 $\Omega$ 的一致正则剖分. 则
\begin{equation}\label{eq:incompressible-hhj-inf-sup}
\sup_{\boldsymbol\tau_h\in\mathcal X_h}
\frac{b_h(\boldsymbol\tau_h,\phi_h)}
{\|\boldsymbol\tau_h\|_{0,h}}
\geq\frac{\beta^*}{\sqrt2C_1}\|\phi_h\|_{2,h}
\qquad\forall\phi_h\in\Phi_h,
\end{equation}
其中 $\beta^*$ 和 $C_1$ 分别为
\eqref{eq:incompressible-sfvp-mesh-inf-sup}
和 \eqref{eq:incompressible-hhj-restriction-stability} 中的常数.
\end{Lemma}

\begin{Remark}\label{rem:incompressible-hhj-w1s-inf-sup}
由 Remark~\ref{rem:incompressible-sfvp-mesh-w1s-inf-sup}, 还有
\[
\sup_{\boldsymbol\tau\in\mathcal X\cap H^1(\Omega)^4}
\frac{b_h(\boldsymbol\tau,\phi_h)}
{\|\boldsymbol\tau\|_{0,h}}
\geq\frac{\gamma(s)}{\sqrt2}|\phi_h|_{1,s,\Omega}
\qquad\forall\text{ real }s\geq2.
\]
所以 Lemma~\ref{lem:incompressible-hhj-inf-sup} 的假设还蕴含
\[
\sup_{\boldsymbol\tau_h\in\mathcal X_h}
\frac{b_h(\boldsymbol\tau_h,\phi_h)}
{\|\boldsymbol\tau_h\|_{0,h}}
\geq\frac{\gamma(s)}{\sqrt2C_1}
|\phi_h|_{1,s,\Omega}
\qquad\forall\text{ real }s\geq2.
\]
\end{Remark}

\begin{Remark}\label{rem:incompressible-hhj-unconditional-uniqueness}
Lemma~\ref{lem:incompressible-hhj-inf-sup} 的构造
还可证明问题 \eqref{eq:incompressible-hhj-discrete-problem}
在不限制 $\Omega$ 和 $\mathcal T_h$ 的情况下有唯一解.
事实上, 因为所用空间有限维, 只需证明
\[
\{\phi_h\in\Phi_h;
\ b_h(\boldsymbol\tau_h,\phi_h)=0
\ \forall\boldsymbol\tau_h\in\mathcal X_h\}
=\{0\}.
\]
如上构造 $\boldsymbol\tau\in\mathcal X\cap H^1(\Omega)^4$, 使
$b_h(\boldsymbol\tau,\phi_h)=|\phi_h|_{1,\Omega}^2$.
由限制正交性又有
$b_h(\Pi_h\boldsymbol\tau,\phi_h)=0$, 故 $\phi_h=0$.
\end{Remark}

最后, 当 $\mathcal T_h$ 一致正则时,
可像 Lemma~\ref{lem:incompressible-sfvp-mesh-zero-norm-equivalence}
一样证明 $\|\cdot\|_{0,h}$ 与 $\|\cdot\|_{0,\Omega}$
是 $\mathcal X_h$ 上一致等价的范数.
证明留作练习, 它来自
\[
\|M_n(\boldsymbol\tau_h)\|_{0,\Gamma_h}
\leq\frac Ch\|\boldsymbol\tau_h\|_{0,\Omega}
\qquad\forall\boldsymbol\tau_h\in\mathcal X_h.
\]
现在可以为问题 \eqref{eq:incompressible-hhj-discrete-problem}
建立最优误差估计.

\MathBookSourcePage{263}
\begin{Theorem}\label{thm:incompressible-hhj-error}
令 $\Omega$ 是有界平面多边形,
$\mathcal T_h$ 是其剖分.
则问题 \eqref{eq:incompressible-hhj-discrete-problem}
有唯一解
$(\boldsymbol\sigma_h,\psi_h)\in\mathcal X_h\times\Phi_h$.
再假设 Stokes 问题 \eqref{eq:incompressible-sfvp-stokes}
的解 $(\mathbf u=\operatorname{curl}\psi,p)$
满足 \eqref{eq:incompressible-hhj-regularity-assumption}, 并令
$\boldsymbol\sigma=(\partial^2\psi/
\partial x_i\partial x_j)_{i,j}$.

\textup{1)} 若 $\mathcal T_h$ 正则, 且
$\boldsymbol\sigma\in H^k(\Omega)^4$, $k\in[1,l]$, 则
\begin{equation}\label{eq:incompressible-hhj-tensor-l2-error}
\|\boldsymbol\sigma-\boldsymbol\sigma_h\|_{0,\Omega}
\leq C_1h^k|\psi|_{k+2,\Omega}.
\end{equation}
若此外 $\Omega$ 为凸域, 则当 $l=1$, $\psi\in H^3(\Omega)$ 时
\begin{equation}\label{eq:incompressible-hhj-linear-h1-error}
|\psi-\psi_h|_{1,\Omega}
\leq C_2h|\psi|_{3,\Omega},
\end{equation}
而当 $l\geq2$, $k\in[2,l]$, $\psi\in H^{k+2}(\Omega)$ 时
\begin{equation}\label{eq:incompressible-hhj-h1-error}
|\psi-\psi_h|_{1,\Omega}
\leq C_3h^k|\psi|_{k+2,\Omega}.
\end{equation}

\textup{2)} 若 $\mathcal T_h$ 一致正则, 则对
$k\in[1,l]$, $\psi\in H^{k+2}(\Omega)$,
\begin{equation}\label{eq:incompressible-hhj-tensor-mesh-error}
\|\boldsymbol\sigma-\boldsymbol\sigma_h\|_{0,h}
\leq C_4h^k|\psi|_{k+2,\Omega}.
\end{equation}
若此外 $\Omega$ 为凸域, $l\geq2$,
$k\in[1,l-1]$, $\psi\in H^{k+2}(\Omega)$, 则
\begin{equation}\label{eq:incompressible-hhj-stream-mesh-error}
\|\psi-\psi_h\|_{2,h}
\leq C_5h^k|\psi|_{k+2,\Omega}.
\end{equation}
\end{Theorem}

\begin{Proof}
像通常一样,
\begin{equation}\label{eq:incompressible-hhj-error-orthogonality}
\left\{
\begin{aligned}
b_h(\boldsymbol\sigma-\boldsymbol\sigma_h,\phi_h)&=0
&&\forall\phi_h\in\Phi_h,\\
a_h(\boldsymbol\sigma-\boldsymbol\sigma_h,\boldsymbol\tau_h)
+b_h(\boldsymbol\tau_h,\psi-\psi_h)&=0
&&\forall\boldsymbol\tau_h\in\mathcal X_h.
\end{aligned}
\right.
\end{equation}
由 Lemma~\ref{lem:incompressible-hhj-restriction-properties}
和 Corollary~\ref{cor:incompressible-hhj-stream-interpolant},
\eqref{eq:incompressible-hhj-error-orthogonality} 给出
\[
b_h(\Pi_h\boldsymbol\sigma-\boldsymbol\sigma_h,\phi_h)=0
\qquad\forall\phi_h\in\Phi_h,
\]
以及
\[
a_h(\boldsymbol\sigma-\boldsymbol\sigma_h,
\Pi_h\boldsymbol\sigma-\boldsymbol\sigma_h)=0.
\]
最后一式直接蕴含
\begin{equation}\label{eq:incompressible-hhj-best-tensor-l2}
\|\boldsymbol\sigma-\boldsymbol\sigma_h\|_{0,\Omega}
\leq\|\boldsymbol\sigma-\Pi_h\boldsymbol\sigma\|_{0,\Omega}.
\end{equation}
当 $\mathcal T_h$ 一致正则时, 两个范数的等价性又给出
\begin{equation}\label{eq:incompressible-hhj-best-tensor-mesh}
\|\boldsymbol\sigma-\boldsymbol\sigma_h\|_{0,h}
\leq C_6\|\boldsymbol\sigma-\Pi_h\boldsymbol\sigma\|_{0,h}.
\end{equation}
于是 \eqref{eq:incompressible-hhj-tensor-l2-error}
和 \eqref{eq:incompressible-hhj-tensor-mesh-error}
由 \eqref{eq:incompressible-hhj-restriction-error} 得到.

其次, \eqref{eq:incompressible-hhj-error-orthogonality}
第二式和 Corollary~\ref{cor:incompressible-hhj-stream-interpolant} 给出
\begin{equation}\label{eq:incompressible-hhj-stream-error-identity}
b_h(\boldsymbol\tau_h,I_h\psi-\psi_h)
=a_h(\boldsymbol\sigma_h-\boldsymbol\sigma,\boldsymbol\tau_h)
\qquad\forall\boldsymbol\tau_h\in\mathcal X_h.
\end{equation}
当 $\Omega$ 为凸域且 $\mathcal T_h$ 一致正则时,
由 Lemma~\ref{lem:incompressible-hhj-inf-sup},
\begin{equation}\label{eq:incompressible-hhj-interpolated-stream-mesh-bound}
\|I_h\psi-\psi_h\|_{2,h}
\leq C_7\|\boldsymbol\sigma_h-\boldsymbol\sigma\|_{0,\Omega}.
\end{equation}
因此 \eqref{eq:incompressible-hhj-stream-mesh-error}
由 \eqref{eq:incompressible-hhj-interpolated-stream-mesh-bound},
\eqref{eq:incompressible-hhj-tensor-l2-error}
和 Corollary~\ref{cor:incompressible-hhj-stream-interpolant} 得到.

\MathBookSourcePage{264}
为证明 \eqref{eq:incompressible-hhj-linear-h1-error}
和 \eqref{eq:incompressible-hhj-h1-error}, 使用熟悉的对偶论证.
对 $\mathbf g\in L^2(\Omega)^2$, 引入辅助 Stokes 问题
\begin{equation}\label{eq:incompressible-hhj-dual-problem}
\left\{
\begin{aligned}
b_h(\boldsymbol\mu_g,\phi)&=(\mathbf g,\operatorname{curl}\phi)
&&\forall\phi\in\widetilde\Psi,\\
a_h(\boldsymbol\mu_g,\boldsymbol\tau)
+b_h(\boldsymbol\tau,\lambda_g)&=0
&&\forall\boldsymbol\tau\in\mathcal X.
\end{aligned}
\right.
\end{equation}
由于 $\Omega$ 为凸域, 解
$(\boldsymbol\mu_g,\lambda_g)$ 属于
$H^1(\Omega)^4\times H^3(\Omega)$, 且
\begin{equation}\label{eq:incompressible-hhj-dual-regularity}
\|\boldsymbol\mu_g\|_{1,\Omega}
+\|\lambda_g\|_{3,\Omega}
\leq C_8\|\mathbf g\|_{0,\Omega}.
\end{equation}
组合
\eqref{eq:incompressible-hhj-error-orthogonality},
\eqref{eq:incompressible-hhj-dual-problem},
\eqref{eq:incompressible-hhj-restriction-orthogonality}
和 Corollary~\ref{cor:incompressible-hhj-stream-interpolant}, 得
\[
\begin{split}
(\mathbf g,\operatorname{curl}(\psi-\psi_h))
={}&b_h(\boldsymbol\mu_g-\Pi_h\boldsymbol\mu_g,
\psi-\phi_h)\\
&+a_h(\boldsymbol\sigma-\boldsymbol\sigma_h,
\boldsymbol\mu_g-\boldsymbol\tau_h)\\
&+b_h(\boldsymbol\sigma-\Pi_h\boldsymbol\sigma,
\lambda_g-I_h\lambda_g)
\end{split}
\]
对所有 $\phi_h\in\Phi_h$,
$\boldsymbol\tau_h\in\mathcal X_h$ 成立.
当 $l\geq2$ 时, Corollary~\ref{cor:incompressible-hhj-stream-interpolant},
\eqref{eq:incompressible-hhj-restriction-error}
和 \eqref{eq:incompressible-hhj-dual-regularity} 给出
\begin{equation}\label{eq:incompressible-hhj-h1-duality-bound}
|\psi-\psi_h|_{1,\Omega}
\leq C_9h\left\{
\inf_{\phi_h\in\Phi_h}\|\psi-\phi_h\|_{2,h}
+\|\boldsymbol\sigma-\boldsymbol\sigma_h\|_{0,\Omega}
+\inf_{\boldsymbol\tau_h\in\mathcal X_h}
\|\boldsymbol\sigma-\boldsymbol\tau_h\|_{0,h}
\right\}.
\end{equation}
当 $l=1$ 时, 只能得到
\begin{equation}\label{eq:incompressible-hhj-linear-h1-duality-bound}
|\psi-\psi_h|_{1,\Omega}
\leq C_{10}h
\left\{
\inf_{\phi_h\in\Phi_h}\|\psi-\phi_h\|_{2,h}
+\|\boldsymbol\sigma-\boldsymbol\sigma_h\|_{0,\Omega}
\right\}
+C_{11}h
\inf_{\boldsymbol\tau_h\in\mathcal X_h}
\|\boldsymbol\sigma-\boldsymbol\tau_h\|_{0,h}.
\end{equation}
利用 \eqref{eq:incompressible-hhj-tensor-l2-error},
Corollary~\ref{cor:incompressible-hhj-stream-interpolant}
和 \eqref{eq:incompressible-hhj-restriction-error},
分别得到
\eqref{eq:incompressible-hhj-h1-error}
和 \eqref{eq:incompressible-hhj-linear-h1-error}.
\end{Proof}

\begin{Corollary}\label{cor:incompressible-hhj-w1s-error}
保留 Theorem~\ref{thm:incompressible-hhj-error} 的全部假设.
若 $l\geq1$, 某个实数 $k\in[1,l]$ 满足
$\psi\in H^{k+2}(\Omega)$, 则对每个 $s>2$,
\begin{equation}\label{eq:incompressible-hhj-w1s-error}
|\psi-\psi_h|_{1,s,\Omega}
\leq C_sh^k|\psi|_{k+2,\Omega}.
\end{equation}
\end{Corollary}

\begin{Proof}
由 \eqref{eq:incompressible-hhj-stream-error-identity}
和 Remark~\ref{rem:incompressible-hhj-w1s-inf-sup},
\begin{equation}\label{eq:incompressible-hhj-interpolated-w1s-bound}
|I_h\psi-\psi_h|_{1,s,\Omega}
\leq C(s)\|\boldsymbol\sigma-\boldsymbol\sigma_h\|_{0,\Omega}.
\end{equation}
结合 \eqref{eq:incompressible-hhj-tensor-l2-error}
和附录插值估计即得结论.
\end{Proof}

\begin{Remark}\label{rem:incompressible-hhj-optimality}
上述 Theorem 显然对所有次数多项式都给出非常整齐的最优误差估计,
而且所考虑的格式相当经济.
另一方面, 所有结果都假设右端项
$\mathbf f\in L^2(\Omega)^2$,
而常常需要求解 $\mathbf f\in L^r(\Omega)^2$, $r<2$ 的 Stokes 问题.
下一小节把误差分析推广到这种情形.
\end{Remark}
